\section{The DC-DFT Lagrangian and the total nuclear gradient}
\label{sec:dc-lagrangian}

\subsection{A Lagrangian with the HF constraint}

Let $C$ denote the HF orbital coefficients and define the VO block of the HF
stationarity residual as
\begin{equation}
  f_{ai}^{\vk}(C,\vR)
  =\left[C_v^{\vk\dagger}F_{\mathrm{HF}}^{\vk}C_o^{\vk}\right]_{ai}.
  \label{eq:hf-stationarity-residual}
\end{equation}
For a converged HF solution, $f=0$.  The target energy is not stationary under
this constraint, so introduce an adjoint multiplier $z$ and define the reduced
Lagrangian
\begin{equation}
  \mathcal L(C,z;\vR)
  =E_{\mathrm{DC-DFA}}(C;\vR)
   +2\operatorname{Re}\sum_{\vk ai}w_\vk
    z_{ai}^{\vk *}f_{ai}^{\vk}(C;\vR).
  \label{eq:dc-lagrangian}
\end{equation}
The multiplier enforcing MO orthonormality is represented below through the
energy-weighted density and the overlap/Pulay term.

Stationarity with respect to orbital rotations,
\begin{equation}
  \frac{\partial\mathcal L}{\partial\kappa}=0,
  \label{eq:lagrangian-orbital-stationarity}
\end{equation}
gives the Z-vector equation of Eq.~\eqref{eq:zvector-equation}.  Explicit
$\dd C/\dd R_{Ax}$ terms then disappear from the total derivative, leaving
\begin{equation}
  \frac{\dd E_{\mathrm{DC-DFA}}}{\dd R_{Ax}}
  =\left.\frac{\partial\mathcal L}{\partial R_{Ax}}\right|_{C,z}
  =g_{Ax}^{\mathrm{explicit}}+g_{Ax}^{\mathrm{response}}.
  \label{eq:dc-gradient-master}
\end{equation}

\subsection{Explicit target contribution}

The explicit derivative of the target energy at fixed density is decomposed as
\begin{equation}
  g_{Ax}^{\mathrm{explicit}}
  =g_{Ax}^{1e}+g_{Ax}^{\mathrm H}+g_{Ax}^{\mathrm{xc}}
   +g_{Ax}^{\mathrm{Pulay,t}}+g_{Ax}^{\mathrm{NN}}.
  \label{eq:explicit-gradient-decomposition}
\end{equation}
The one-electron term is
\begin{equation}
  g_{Ax}^{1e}
  =\frac{1}{N_k}\sum_\vk
    \Tr[D^{\vk}h^{\vk,Ax}],
  \label{eq:one-electron-gradient}
\end{equation}
where $h^{Ax}$ contains the AO-matrix derivatives of the kinetic, local GTH,
and nonlocal GTH operators.

At fixed density, the Hartree term is
\begin{equation}
  g_{Ax}^{\mathrm H}
  =\frac{1}{2N_k}\sum_\vk
    \Tr[D^{\vk}J^{\vk,Ax}[D]].
  \label{eq:hartree-gradient}
\end{equation}
For FFTDF, $J^{Ax}[D]$ contains the AO/grid derivative and the appropriate
Coulomb-kernel convention.  For GDF and RSDF it contains the three-center
source and two-center metric derivatives.

The explicit semilocal XC derivative differentiates the AOs and the
quadrature representation while holding the density matrix fixed,
\begin{equation}
  g_{Ax}^{\mathrm{xc}}
  =\left.
    \frac{\partial E_{\mathrm{xc}}[D;\vR]}{\partial R_{Ax}}
   \right|_D.
  \label{eq:xc-explicit-gradient}
\end{equation}
Atom-centered Gaussian AOs are evaluated on a uniform periodic grid, so the AO
center derivative and periodic-image/Bloch-phase conventions must match those
of the central value.  SCAN and r2SCAN add the $\tau$ channel.

\subsection{The target overlap/Pulay term}

Define the target energy-weighted density from the occupied--occupied block of
the target Fock matrix,
\begin{equation}
  W_{\mathrm t}^{\vk}
  =2C_o^{\vk}
    \left(C_o^{\vk\dagger}F_{\mathrm t}^{\vk}C_o^{\vk}\right)
    C_o^{\vk\dagger}.
  \label{eq:target-energy-weighted-density}
\end{equation}
The target overlap contribution is then
\begin{equation}
  g_{Ax}^{\mathrm{Pulay,t}}
  =-\frac{1}{N_k}\sum_\vk
    \Tr[W_{\mathrm t}^{\vk}S^{\vk,Ax}].
  \label{eq:target-overlap-pulay}
\end{equation}
An AO-row derivative formulation may redistribute parts of
Eqs.~\eqref{eq:one-electron-gradient}--\eqref{eq:target-overlap-pulay} among
atom-row contractions, but their sum must be identical.

Using GTH ionic charges and the periodic Ewald convention, the ion--ion term is
\begin{equation}
  g_{Ax}^{\mathrm{NN}}
  =\frac{\partial E_{\mathrm{NN}}}{\partial R_{Ax}}.
  \label{eq:nuclear-repulsion-gradient}
\end{equation}

\subsection{The HF orbital-response contribution}

Writing Eq.~\eqref{eq:zvector-contraction} explicitly in indices gives
\begin{equation}
  g_{Ax}^{\mathrm{response}}
  =2\operatorname{Re}\sum_{\vk ai}w_\vk
   z_{ai}^{\vk *}B_{ai}^{\vk,Ax}.
  \label{eq:response-gradient}
\end{equation}
Here $B^{Ax}$ is the nuclear derivative of the source HF equation.  Changing
the target between PBE and an MGGA does not change the physical definitions of
$B^{Ax}$ or the HF Hessian.  Only the target VO gradient $g$ and the explicit
XC derivative change, so the same response infrastructure can be reused.

\subsection{Backend-independent assembly}

The final package decomposition is recorded as
\begin{align}
  g_{Ax}^{\mathrm{elec,explicit}}
  &=g_{Ax}^{1e}+g_{Ax}^{\mathrm H}+g_{Ax}^{\mathrm{xc}}
    +g_{Ax}^{\mathrm{Pulay,t}},\\
  g_{Ax}
  &=g_{Ax}^{\mathrm{elec,explicit}}
    +g_{Ax}^{\mathrm{NN}}+g_{Ax}^{\mathrm{response}}.
  \label{eq:final-gradient-assembly}
\end{align}
FFTDF, GDF, and RSDF differ in their providers for J/K values and derivatives;
the Z-vector equation and final assembly contract are shared.

\begin{remark}
A full-energy finite difference reconverges the source HF solution and
re-evaluates the target energy at both $\vR+h\mathbf e_{Ax}$ and
$\vR-h\mathbf e_{Ax}$.  A component finite difference with the central density
held fixed is useful for diagnosing an individual kernel such as
Eq.~\eqref{eq:xc-explicit-gradient}, but it is not a reference for the total
analytic gradient.
\end{remark}

\subsection{Translational invariance}

Translating all atoms by the same vector within a fixed lattice must not change
the energy.  In a converged representation,
\begin{equation}
  \sum_A\mathbf g_A=\mathbf0.
  \label{eq:translation-sum-rule}
\end{equation}
A large net gradient is a sensitive diagnostic for missing AO-image
derivatives, auxiliary-center response, XC-grid aliasing, $G=0$ treatment, or
k-point normalization.
